% =============================================================================
% AI Tageszeitung — Kommune-Startseite
% Zeigt alle aktuellen Nachrichten EINER Gemeinde
% Variable: \kommunename, \kommunefarbe
% =============================================================================
\documentclass[11pt, a4paper]{scrartcl}

\usepackage{fontspec}
\usepackage[ngerman]{babel}
\usepackage[top=12mm, bottom=15mm, left=10mm, right=10mm]{geometry}
\usepackage{multicol}
\usepackage{graphicx}
\usepackage{xcolor}
\usepackage{hyperref}
\usepackage{tikz}
\usepackage{tcolorbox}
\usepackage{datetime2}
\usepackage{longtable}
\usepackage{booktabs}

% --- Farben ---
\definecolor{rbkblau}{HTML}{1B3A5C}
\definecolor{rbkrot}{HTML}{C0392B}
\definecolor{rbkgrau}{HTML}{F5F5F5}
\definecolor{akzent}{HTML}{2980B9}
\definecolor{liniengrau}{HTML}{BFBFBF}

% --- Kommune-spezifisch (wird per Script gesetzt) ---
\newcommand{\kommunename}{%%KOMMUNE_NAME%%}
\newcommand{\kommuneslug}{%%KOMMUNE_SLUG%%}

% --- Schriften ---
\setmainfont{TeX Gyre Termes}
\setsansfont{TeX Gyre Heros}
\newfontfamily\headlinefont{TeX Gyre Heros Bold}

\hypersetup{
  colorlinks=true,
  linkcolor=rbkblau,
  urlcolor=akzent,
}

% --- Artikel-Makro ---
\newcommand{\artikel}[5]{%
  % #1=Datum/Zeit, #2=Kategorie, #3=Headline, #4=Teaser, #5=Link
  \noindent
  \begin{minipage}{\linewidth}
    {\footnotesize\sffamily\color{rbkblau}#1\hfill\textit{#2}}\\[2pt]
    {\large\headlinefont #3}\\[3pt]
    {\small #4}\\[2pt]
    {\footnotesize\color{akzent}\href{#5}{Vollständiger Artikel \textrightarrow}}
  \end{minipage}
  \\[1mm]
  \rule{\linewidth}{0.2pt}\\[3mm]
}

% --- Dokument ---
\begin{document}
\pagestyle{empty}

% === KOPFZEILE ===
\begin{center}
\begin{tikzpicture}
  \fill[rbkblau] (0,0) rectangle (\textwidth, 1.5cm);
  \node[anchor=west, white] at (0.3, 0.75) {%
    {\fontsize{22}{26}\selectfont\headlinefont \kommunename}};
  \node[anchor=east, white] at (\textwidth-0.3, 1.0) {%
    {\small\sffamily AI Tageszeitung}};
  \node[anchor=east, white] at (\textwidth-0.3, 0.4) {%
    {\footnotesize\sffamily Stand: \DTMnow}};
\end{tikzpicture}
\end{center}

\vspace{2mm}

% === NAVIGATION ===
{\small\sffamily
\href{../index.html}{\textbf{Startseite}}~|~
\href{../bergisch-gladbach/index.html}{GL}~|~
\href{../burscheid/index.html}{BUR}~|~
\href{../kuerten/index.html}{KÜR}~|~
\href{../leichlingen/index.html}{LEI}~|~
\href{../odenthal/index.html}{ODE}~|~
\href{../overath/index.html}{OVE}~|~
\href{../roesrath/index.html}{RÖS}~|~
\href{../wermelskirchen/index.html}{WER}
}

\vspace{4mm}
\rule{\textwidth}{0.8pt}
\vspace{4mm}

% === ARTIKEL-LISTE ===
\begin{multicols}{2}
\raggedcolumns

%%ARTIKEL_PLACEHOLDER%%

% Beispiel:
\artikel{05.05.2026, 08:30}{Politik}
  {Ratssitzung beschließt neuen Haushalt}
  {Mit großer Mehrheit wurde der Haushalt 2026/27 verabschiedet. Schwerpunkte: Schulsanierung und Radwegebau.}
  {https://example.com/artikel/detail/42}

\artikel{05.05.2026, 07:15}{Blaulicht}
  {Unfall auf der L286: Zwei Verletzte}
  {Heute Morgen kollidierten zwei PKW auf der Landstraße zwischen Dürscheid und Bechen.}
  {https://example.com/artikel/detail/43}

\end{multicols}

% === FUßZEILE ===
\vfill
\begin{center}
{\footnotesize\sffamily\color{rbkblau}
  AI Tageszeitung~|~Rheinisch-Bergischer Kreis~|~\kommunename\\[1mm]
  \href{../index.html}{Zurück zur Übersicht}
}
\end{center}

\end{document}
